
\scanpage{51}
\mspage{\deleted{\unsure{6}}7}
d'une structure déjà orientée si $a$ est une arête libre (correspondant à un $C\in\pi_0(\partial X)$), ou qu'on dispose de l'une des deux orientations $\omega(a)$ --- dans le cas où $a$ est une arête ordinaire, implique que l'on tient en mains, pour chaque $a\in A$, un polygone combinatoire $P_a$, dont l'ens des sommets est $R_a$, et l'ens des arêtes peut se décrire \deleted{\unsure{aussi}} de façon canonique en termes de $R_a,\omega_a$ et $\omega_a\longrightarrow\operatorname{Circ}(R_a)$.
Il peut être plus commode dans la donnée d'ensemble de la structure, \unsure{où} la famille des $R_a$ serait unifiée par un
\deleted{qui serait \unsure{les uns un contour}}
\uline{contour combinatoire}, i.e. un ensemble de repères $\vec{\underline R}$ (plutôt ``repères de rigidification''), avec une opération de
\[
 \mathbb D_\infty=\{\sigma_0,\sigma_1\mid\sigma_0^2=\sigma_1^2=1\}
\]
dessus, et une application
\[
 \vec{\underline R}\longrightarrow\vec A
\]
compatible avec l'hom.
\[
 \mathbb D_\infty\xrightarrow{\varepsilon}\{\pm1\}
\]
sur les groupes d'opérateurs, mais de plus un sous ens
\[
 \vec{\underline R}_l^+\subset\vec{\underline R}\mid\vec A_l
\]
définissant une orientation du contour correspondant.
\margin{$(\vec A_l=\text{arêtes libres, correspondant aux cercles de recollements, i.e. les }C\in\pi_0(\unsure{\partial}X))$}
\mspage{50}
